\title{xLSTM: Extended Long Short-Term Memory}

\begin{abstract}
The constant error carousel and gating mechanisms, foundational to Long Short-Term Memory (LSTM), were developed in the 1990s. LSTMs have since demonstrated enduring relevance and have powered many advances in deep learning, notably forming the basis for the earliest Large Language Models (LLMs). Nevertheless, with the emergence of Transformer models—centered around parallelizable self-attention—LSTMs were largely superseded in terms of scalability and performance. This leads us to a straightforward inquiry: To what extent can language modeling be improved by scaling LSTMs to billions of parameters, utilizing modern LLM methodologies, while addressing LSTM-specific limitations? To answer this, we propose two main innovations: introducing exponential gating coupled with tailored normalization and stabilization strategies; and reconfiguring the LSTM memory architecture to create (i) sLSTM, characterized by a scalar memory, scalar update mechanism, and revised memory mixing, and (ii) mLSTM, featuring a matrix-based memory and a covariance-driven update rule, enabling full parallelism. When these augmented LSTM designs are incorporated within residual block frameworks, the resulting xLSTM blocks can be stacked to construct xLSTM-based architectures. The implementation of exponential gating and reimagined memory structures equips xLSTMs to rival leading Transformers and State Space Models, both in terms of efficiency and scalability.\\
Code available at: \url{https://github.com/NX-AI/xlstm}
\end{abstract}

\section{Introduction}
\label{sec:intro}

The concepts behind Long Short-Term Memory (LSTM) networks~\citep{Hochreiter:91,Hochreiter:97nips,Hochreiter:97}, particularly the constant error carousel and the use of gating mechanisms, were proposed as solutions to the vanishing gradient issue encountered in recurrent neural networks~\citep{Hochreiter:91,Hochreiter:00book}:

\begin{align}
\label{eq:lstm_idea}
  c_t \ = \ f_t \ c_{t-1} \ + \ 
          i_t \  z_t \ , \quad  
          h_t \ = \ o_t \ \psi({c_t}) \ . 
\end{align}

The constant error carousel operates by summing the previous cell state $c_{t-1}$ (shown in green) with new cell inputs $z_t$, where this process is regulated by sigmoid gates (depicted in blue). The input gate $i_t$ and forget gate $f_t$ determine the dynamics of this updating mechanism, whereas the output gate $o_t$ dictates how the memory cell's value is translated into the hidden state $h_t$. The cell state is first passed through a squashing function $\psi$ for normalization, after which the output gate applies gating to produce the hidden state.

Long Short-Term Memory networks (LSTMs) have found widespread adoption across a range of domains \citep{Hochreiter:01,Hochreiter:07,Schmidhuber:15}, maintaining prominence in text generation prior to the rise of Transformers in 2017~\citep{Vaswani:17}. Their utility has been substantiated through diverse sequence processing tasks, including but not limited to text generation~\citep{Graves:13,Karpathy:15blog}, handwriting synthesis~\citep{Graves:13}, sequence-to-sequence translation~\citep{Sutskever:14nips}, program evaluation~\citep{Zaremba:14arxiv}, image captioning~\citep{Karpathy:15,Hossain:19}, code generation~\citep{Karpathy:15blog}, rainfall-runoff simulations~\citep{Kratzert:18,Kratzert:19}, and hydrological modeling for flood alerts~\citep{Nearing:24}. In reinforcement learning, LSTMs stand out as the top-performing sequence architectures, as exemplified by their deployment in AlphaStar for StarCraft II~\citep{Vinyals:17short}, OpenAI Five for Dota 2~\citep{OpenAI:19}, and models for magnetic control in nuclear fusion experiments~\citep{Degrave:22short}. Notably, LSTMs demonstrate proficiency in learning high-level abstractions—effectively capturing semantic information within their memory structures~\citep{Karpathy:15blog}. This capacity is illustrated by discoveries such as neurons encoding numerical and syntactic structure~\citep{Lakretz:19}, linguistic features~\citep{Bau:19}, or sentiment~\citep{Radford:17arxiv}. Despite emerging alternatives, LSTMs continue to underpin critical real-world applications~\citep{Degrave:22short,Nearing:24}, underscoring their enduring significance.

Although LSTMs have achieved significant advancements, they exhibit three primary shortcomings: (i) They lack the ability to update previously made storage decisions. This issue is illustrated by the {\it Nearest Neighbor Search} task (see also Appendix~\ref{sec:appNearestNeighborSearch}): Here, a model is required to traverse a sequence to find the vector most similar to a reference and retrieve its corresponding value at the end of the sequence. As depicted in the left panel of Figure~\ref{fig:lstmProblems}, LSTM models display higher mean squared error on this task because they are not adept at overwriting stored information when a better match appears later in the sequence. Our proposed xLSTM, utilizing exponential gating, addresses this deficit. (ii) They possess restricted storage capability, as all relevant information must be encoded into one-dimensional cell states. We demonstrate this through the {\it Rare Token Prediction} task; the right panel of Figure~\ref{fig:lstmProblems} reports the perplexity scores for Wikitext-103~\citep{Merity:17} tokens separated by their frequency. The LSTM shows inferior performance for infrequent tokens due to its storage constraints, while our xLSTM, equipped with a matrix-based memory, overcomes this barrier. (iii) They are not easily parallelized because memory content from one step is mixed with that from the next through hidden-to-hidden connections, necessitating strictly sequential computation.

The shortcomings of LSTM have led to the rise of Transformers~\citep{Vaswani:17} within the field of language modeling. This raises the question: if these LSTM limitations are addressed and LSTMs are scaled up to match the parameter count of modern Large Language Models, what level of performance is attainable in language modeling?

\section{Extended Long Short-Term Memory}
\label{sec:xLSTM}

To address the shortcomings inherent in the LSTM architecture, Extended Long Short-Term Memory (xLSTM) proposes two primary enhancements to the foundational LSTM concept as described in Equation~\eqref{eq:lstm_idea}. The xLSTM expands the LSTM family through the incorporation of (i) exponential gating and (ii) innovative memory architectures, thus leading to two novel variants: the sLSTM (refer to Section~\ref{sec:sLSTM}), characterized by its scalar memory, scalar update mechanism, and use of memory mixing; and the mLSTM (discussed in Section~\ref{sec:mLSTM}), which operates with a matrix-based memory and employs an update rule grounded in covariance (outer product), enabling full parallelization. Both these variants augment the traditional LSTM using exponential gating schemes. In particular, the mLSTM achieves parallelizability by dispensing with memory mixing, effectively eliminating hidden-to-hidden recurrent links. Extensions to multiple memory cells are possible for both the sLSTM and mLSTM frameworks, with sLSTM specifically allowing memory mixing among its cells. Moreover, sLSTM may include multiple heads, whereby memory mixing occurs solely across cells within each individual head and not between heads themselves. This architectural extension—coupling the concept of heads with exponential gating—introduces a new modality of memory mixing in sLSTM. For mLSTM, the distinction between having multiple heads and multiple cells collapses, as these two configurations are functionally equivalent.

By embedding these novel LSTM variants within residual block modules, we construct xLSTM blocks (refer to Section~\ref{sec:xLSTMblock}). When these xLSTM blocks are assembled in a residual fashion across architectural layers, they form what we term xLSTM architectures (see Section~\ref{sec:xLSTMarchitecure}). The structure and constituent parts of the xLSTM architecture are illustrated in Figure~\ref{fig:xlstm_sketch}.

\subsection{Review of the Long Short-Term Memory}
\label{sec:orgLSTM}

The initial conception of the LSTM architecture~\citep{Hochreiter:91,Hochreiter:97nips,Hochreiter:97} featured a scalar memory cell serving as the fundamental component for information processing and retention. This design addresses the vanishing gradient problem~\citep{Hochreiter:91,Hochreiter:00book} by employing the constant error carousel mechanism, realized via the cell state update. Three gating mechanisms—namely, the input, output, and forget gates—regulate the flow of information within the cell. The forget gate was subsequently incorporated into the architecture by \citet{Gers:00}. The update equations governing the behavior of the LSTM memory cell at each time step $t$ are given by:

\begin{align}
c_t \ &= \  f_t \ c_{t-1} \ + \ i_t \ z_t &  & &\text{cell state} \\
h_t  \ &= \ o_t \ \tilde{h}_t \ , & \tilde{h}_t \ &= \ \psi \left( c_t\right)
  &\text{hidden state} \\
z_t \ &= \ \varphi \left( \tilde{z}_t \right) \ , 
  &\tilde{z}_t \ &=  \ \mathbf{w}^\top_{z} \ \mathbf{x}_t \ + \
  r_{z}  h_{t-1} \ + \  b_{z} \ \
  &\text{cell input} \\
i_t \ &= \ \sigma \left( \tilde{i}_t  \right) \ , 
  &\tilde{i}_t \ &= \ \mathbf{w}^\top_{i} \ \mathbf{x}_t \ + \
  r_{i}  \ h_{t-1} \ + \  b_{i} \ \
  &\text{input gate} \\
f_t \ &= \ \sigma \left(  \tilde{f}_t \right) \ , 
  &\tilde{f}_t \ &= \ \mathbf{w}^\top_{f} \ \mathbf{x}_t  \ + \
  r_{f}  \ h_{t-1} \ + \  b_{f} \ \
  &\text{forget gate} \\
o_t \ &= \ \sigma \left( \tilde{o}_t \right) \ , 
  &\tilde{o}_t  \ &= \ \mathbf{w}^\top_{o} \ \mathbf{x}_t \ + \
  r_{o}  \ h_{t-1} \ + \  b_{o} \ \
  &\text{output gate} 
\end{align}

The vectors $\mathbf{w}_{z}$, $\mathbf{w}_{i}$, $\mathbf{w}_{f}$, and $\mathbf{w}_{o}$ represent the weights connecting the inputs $\mathbf{x}_t$ to the cell input, input gate, forget gate, and output gate, respectively. The parameters $r_{z}$, $r_{i}$, $r_{f}$, and $r_{o}$ denote the recurrent weights linking the previous hidden state $h_{t-1}$ to each of the cell input, input gate, forget gate, and output gate, in that order. The bias terms associated with each component are denoted as $b_{z}$, $b_{i}$, $b_{f}$, and $b_{o}$. The activation functions for the cell input and hidden state, denoted as $\varphi$ and $\psi$ respectively, are usually instantiated as $\tanh$. The function $\psi$ serves to constrain the cell state, which would otherwise not be limited in range. All gates employ a sigmoid nonlinearity, expressed as $\sigma \left( x \right)= 1/(1+\exp(-x))$. 

Subsequent developments introduced a vector formulation for the memory cell, $\mathbf{c}_t \in \mathbb{R}^{d}$, consolidating several scalar cells and enabling the application of recurrent weight matrices $\mathbf{R} \in \mathbb{R}^{d \times d}$ that mix outputs across memory cells~\citep{Greff:15}; additional details are available in Appendix~\ref{sec:apporgLSTM}. Empirical ablation studies have demonstrated the necessity of all memory cell components for effective operation~\citep{Greff:15}.

\subsection{sLSTM}
\label{sec:sLSTM}

In order to enhance LSTMs with mechanisms for adjusting their storage choices, we propose the inclusion of exponential gates (shown in red), accompanied by techniques for normalization and stabilization. Specifically, both the input and forget gates may utilize exponential activation functions. For normalization, we define a normalizer state, which accumulates the sum of the input gate multiplied by all subsequent forget gates.

The scalar sLSTM forward propagation proceeds as follows:

\begin{align}
\label{eq:slstmforward}
c_t \ &= \  f_t \ c_{t-1} \ + \ i_t \ z_t & & 
 &\text{cell state} \\
n_t \ &= \  f_t \ n_{t-1} \ + \ i_t  & & 
 &\text{normalizer state} \\
h_t \ &= \ o_t \ \tilde{h}_t \ ,
  &\tilde{h}_t \ &= \ c_t / n_t
&\text{hidden state} \\
z_t \ &= \ \varphi \left( \tilde{z}_t \right) \ , 
  &\tilde{z}_t \ &=  \ \mathbf{w}^\top_{z} \ \mathbf{x}_t \ + \
  r_{z}  h_{t-1} \ + \  b_{z}
  &\text{cell input} \\
i_t \ &= \ \!\exp \left( \tilde{i}_t  \right) \ , 
  &\tilde{i}_t \ &= \ \mathbf{w}^\top_{i} \ \mathbf{x}_t \ + \
  r_{i}  \ h_{t-1} \ + \  b_{i}
  &\text{input gate} \\
f_t \ &= \ \sigma \left(  \tilde{f}_t \right) \ \text{OR} \
\!\exp \left(  \tilde{f}_t \right) \ ,
  &\tilde{f}_t \ &= \ \mathbf{w}^\top_{f} \ \mathbf{x}_t  \ + \
  r_{f}  \ h_{t-1} \ + \  b_{f}
  &\text{forget gate} \\
o_t \ &= \ \sigma \left( \tilde{o}_t \right) \ , 
  &\tilde{o}_t  \ &= \ \mathbf{w}^\top_{o} \ \mathbf{x}_t \ + \
  r_{o}  \ h_{t-1} \ + \  b_{o}
  &\text{output gate} 
\end{align}

We adapt the traditional LSTM gating strategies—specifically, those involving input- and hidden-state dependent gates along with bias terms—to these alternative architectures. Given that exponential activation functions have the potential to generate excessively large outputs and result in overflows, we incorporate an extra state \mbox{$m_t$~\citep{Milakov:18arxiv}} to stabilize the gates:

\begin{align}
\label{eq:slstmstabil}
m_t \ &= \ \max \left( \log ( f_t ) + m_{t-1} , \log ( i_t ) \right) &\text{stabilizer state} \\
i'_t \ &= \ \!\exp \left( \log \left ( i_t \right) - m_t \right) \ = \ \!\exp \left( \tilde{i}_t - m_t \right) \ \, 
  &\text{stabil. input gate} \\
f'_t \ &= \ \!\exp \left( \log \left( f_t \right) + m_{t-1} - m_t \right) \ \, 
  &\text{stabil. forget gate}
\end{align}

As demonstrated in Appendix~\ref{sec:appsLSTM}, substituting $f_t$ with $f'_t$ and $i_t$ with $i'_t$ during the forward pass leaves both the overall network output and the loss gradients with respect to the parameters unaffected.

\paragraph{New Memory Mixing.} Similar to the original LSTM, sLSTM supports the use of several memory cells (refer to Appendix~\ref{sec:appsLSTM}). By leveraging recurrent connections—$\mathbf{R}_{\mathbf{z}}$, $\mathbf{R}_{\mathbf{i}}$, $\mathbf{R}_{\mathbf{f}}$, and $\mathbf{R}_{\mathbf{o}}$—that link the hidden state $\mathbf{h}$ to both the memory cell input $\mathbf{z}$ and the corresponding gates $\mathbf{i}$, $\mathbf{f}$, and $\mathbf{o}$, multiple memory cells facilitate mixing of information. What distinguishes the updated memory mixing mechanism is the application of exponential gating. The re-envisioned sLSTM architecture incorporates several heads, allowing for memory mixing within a single head but not between different heads. Thus, the combination of sLSTM heads with exponential gating introduces a novel framework for memory mixing.

\subsection{mLSTM}
\label{sec:mLSTM}

In order to bolster the storage capabilities of LSTMs, we generalize the memory cell from a single scalar $c \in \mathbb{R}$ to a full matrix $\mathbf{C} \in \mathbb{R}^{d \times d}$. Retrieval operations are consequently achieved using matrix multiplication. Specifically, at each time step $t$, the objective is to store a key-value pair: the key $\mathbf{k}_t \in \mathbb{R}^d$ and the value $\mathbf{v}_t \in \mathbb{R}^d$, adopting the naming conventions from the Transformer literature. Subsequently, at some future time $t + \tau$, the stored value $\mathbf{v}_t$ should be recalled using a query vector $\mathbf{q}_{t+\tau} \in \mathbb{R}^d$. This framework corresponds to the paradigm of Bidirectional Associative Memories (BAMs) \citep{Kohonen:72,Anderson:72,Nakano:72,Anderson:77}.

To encode the key-value pair, we utilize the covariance update rule~\citep{Sejnowski:77,Dayan:91}

\begin{align}
\mathbf{C}_t \ &= \ \mathbf{C}_{t-1} \ + \ \mathbf{v}_{t} \ \mathbf{k}_t^\top \ .
\end{align}

We consider that inputs are normalized with layer-norm prior to their projection as keys and values, ensuring a mean of zero. The covariance update mechanism is proven to be optimal~\citep{Dayan:91} for achieving maximal separability among retrieved binary vectors, effectively maximizing the signal-to-noise ratio. Enhanced separability can be realized if retrieval is restricted to pairwise terms and a quadratic computational cost—as in attention mechanisms—is accepted~\citep{Krotov:16,Krotov:17,Ramsauer:21}. Notably, this covariance update is identical to the approach used in Fast Weight Programmers \citep{Schmidhuber:92ncfastweights,Schlag:21}, which later introduced a fixed decay parameter for $\mathbf{C}_{t-1}$ and a fixed learning parameter for $\mathbf{v}_{t} \mathbf{k}_t^\top$~\citep{Ba:16}. Following this idea, we adapt the covariance update to the LSTM architecture, mapping the forget gate to the decay rate, the input gate to the learning rate, and the output gate to the scaling of the retrieved output.

Within this matrix memory framework, the normalizer state represents a sum of key vectors, with each key's contribution scaled by the input gate along with the product of subsequent forget gates. This construction ensures that the normalizer encodes information about the gate intensities over time. Because the dot product between a query and the normalizer state might approach zero, we employ the absolute value of the dot product and enforce a minimum threshold (usually set at 1.0), consistent with the method in~\citep{Sun:23arxiv}. The forward computation for mLSTM proceeds as follows:

\begin{align}
\label{eq:mlstm_recurrent_begin}
\mathbf{C}_t \ &= \  f_t \ \mathbf{C}_{t-1} \ + \ 
  i_t \ \mathbf{v}_t \ \mathbf{k}_t^\top & &  &\text{cell state} \\
\mathbf{n}_t \ &= \  f_t \ \mathbf{n}_{t-1} \ + \ 
  i_t \ \mathbf{k}_t & &  &\text{normalizer state} \\
\mathbf{h}_t  \ &= \ \mathbf{o}_t \ \odot \ \tilde{\mathbf{h}}_t
\ , \qquad \qquad 
\tilde{\mathbf{h}}_t \ = \   \mathbf{C}_t \mathbf{q}_t \ / \ 
  \max \left\{ |\mathbf{n}_t^\top \mathbf{q}_t|, 1 \right\} 
   &&&\text{hidden state} \\
\mathbf{q}_t \ &= \ \mathbf{W}_q \ \mathbf{x}_t \ + \ \mathbf{b}_q  & &  &\text{query input} \\
\mathbf{k}_t \ &= \ \frac{1}{\sqrt{d}} \mathbf{W}_k \ \mathbf{x}_t \ + \ \mathbf{b}_k  & &  &\text{key input} \\
\mathbf{v}_t \ &= \ \mathbf{W}_v \ \mathbf{x}_t \ + \ \mathbf{b}_v  & &  &\text{value input} \\
i_t \ &= \ \!\exp \!  \! \left( \tilde{i}_t  \right) \ , 
 \qquad \qquad \ \ \ \ \,
  \tilde{i}_t \ = \ \mathbf{w}^\top_{i} \ \mathbf{x}_t \ + \  b_{i} \ \ \ 
  &&&\text{input gate} \\
f_t \ &= \ \sigma \!  \left(  \tilde{f}_t \right) \ \text{OR} \ \!\exp \!  \! \left(  \tilde{f}_t \right) , \ \ \: \!
\tilde{f}_t \ = \ \mathbf{w}^\top_{f} \ \mathbf{x}_t  \ + \
  b_{f}
  &&& \text{forget gate} \\
\label{eq:mlstm_recurrent_end}
\mathbf{o}_t \ &= \ \sigma \left( \tilde{\mathbf{o}}_t \right) \ , \qquad \qquad \quad \ \ \ \,
  \tilde{\mathbf{o}}_t  \ = \ \mathbf{W}_{\mathbf{o}} \ \mathbf{x}_t \ + \
  \mathbf{b}_{\mathbf{o}}
  &&&\text{output gate} 
\end{align}

mLSTM is capable of supporting several memory cells, similar to the standard LSTM architecture. In the context of mLSTM, having several heads is functionally the same as having multiple cells, owing to the absence of any memory mixing between them. To address the instability that can arise from the exponential gates in mLSTM, we apply the same stabilization methods utilized for sLSTM (refer to Equation~\ref{eq:slstmstabil}).  
Because memory states do not interact within mLSTM, the recurrence relation can be restated in a form that allows for parallel computation. Additional information is provided in Appendix~\ref{sec:appmLSTM}.

\subsection{xLSTM Architecture}
\label{sec:xLSTMarch}

\paragraph{xLSTM Blocks.}
\label{sec:xLSTMblock}
An xLSTM block is designed to encode previous information non-linearly within a high-dimensional representation, thereby enhancing the separation among distinct contexts or historical sequences. The ability to distinguish between different histories is essential for accurate prediction of subsequent elements in a sequence, such as the next token. This approach is inspired by Cover’s Theorem~\citep{Cover:65}, which asserts that non-linear embeddings into higher-dimensional spaces increase the likelihood that different patterns can be linearly separated, as opposed to the separation achievable in their original dimensionality. We analyze two types of residual block structures: (i) A residual block with an up-projection occurring after the non-linear transformation (akin to what is found in Transformers), wherein the historical sequence is non-linearly processed within the original space, followed by a linear mapping to a higher-dimensional space, a subsequent non-linear activation, and finally a linear mapping returning to the original dimension; refer to the left side of Figure~\ref{fig:Backbones} and the third column of Figure~\ref{fig:xlstm_sketch}. An expanded depiction can be found in Figure~\ref{fig:appsLSTM_detailed} in the appendix. (ii) A residual block with an up-projection preceding the non-linearity (similar to State Space Models), in which the input is first mapped linearly to a higher-dimensional space,

It performs a nonlinear aggregation of previous information within a high-dimensional representation, followed by a linear transformation back to the initial space. When an xLSTM block includes an sLSTM, we apply the post up-projection block; conversely, for an xLSTM block with an mLSTM, we opt for the pre up-projection block due to the increased memory capacity provided by the higher-dimensional setting. For further clarification, see the left panel of Figure~\ref{fig:Backbones}, the third column of Figure~\ref{fig:xlstm_sketch}, or consult Figure~\ref{fig:appsLSTM_detailed} in the appendix.

\paragraph{xLSTM Architecture. }
\label{sec:xLSTMarchitecure}
The xLSTM model is formed by stacking multiple modular components in a residual fashion~\citep{Srivastava:15,He:16}. This design leverages the widely adopted pre-LayerNorm~\citep{Ba:16b} residual structure that serves as the backbone in state-of-the-art Large Language Models. Refer to the final column of Figure~\ref{fig:xlstm_sketch} for an illustration.

\subsection{Memory and Speed Considerations}

Unlike Transformers, xLSTM networks exhibit linear computational complexity and maintain a fixed memory requirement regardless of the length of the sequence. Owing to the compressive nature of xLSTM memory, these networks are particularly advantageous for deployment in industrial contexts and for edge computing scenarios.

While the memory component in mLSTM is parameter-free, its operations involve intensive computations due to the use of a $d \times d$ memory matrix and $d \times d$ updates. Our approach balances the need for memory capacity with the associated computational load. Despite this computational intensity, these operations are amenable to parallelization on GPUs, resulting in only a negligible impact on actual run time.

Although mLSTM can be parallelized similarly to methods such as FlashAttention~\citep{Dao:22,Dao:23} or GLA~\citep{Yang:23arxiv}, sLSTM lacks parallelizability because of the memory mixing that involves hidden-to-hidden connections. Nevertheless, we introduced an efficient CUDA-based implementation that leverages GPU memory optimization down to the register level, resulting in performance that is generally no more than twice as slow as that of mLSTM.

\section{Related Work}
\label{sec:related}

\paragraph{Linear Attention.} A variety of techniques have been proposed to address the Transformer’s quadratic scaling with respect to sequence length by achieving linear computational complexity in attention. For instance, the Synthesizer dispenses with direct token-to-token dependencies by learning artificial attention patterns~\citep{Tay:20arxiv}. The Linformer approximates self-attention using low-rank matrices, thus enabling linear computation~\citep{Wang:20arxiv}. The Linear Transformer modifies the standard attention computation to make it linear~\citep{Katharopoulos:20}. The Performer leverages positive orthogonal random features to approximate the softmax attention mechanism in a linear manner~\citep{Choromanski:21}. Moreover, attention operations have been supplanted by rapid long-range convolutional layers in models such as Structured Global Convolution (SGConv)~\citep{Li:22} and the Hyena Hierarchy~\citep{Poli:23}.

\paragraph{State Space Models.} In recent years, State Space Models (SSMs) have attracted significant attention due to their linear scaling with sequence length and competitive results against Transformers. Initial advancements in this area began with the introduction of the Structured State Space sequence model (S4)~\citep{Gu:21}, which was soon succeeded by subsequent models including the Diagonal State Space (DSS) model~\citep{Gupta:22}, Gated State Space (GSS) models~\citep{Mehta:22}, S5 model~\citep{Smith:22}, Bidirectional Gated SSM (BiGS)~\citep{Wang:22}, H3 model~\citep{Fu:23}, and Mamba~\citep{Gu:24arxiv}.

\paragraph{Recurrent Neural Networks.} Recurrent Neural Networks (RNNs) have been proposed as alternatives to Transformers and attention mechanisms, mainly due to their favorable linear scaling with respect to context length. Recent advancements, such as RNNs employing Deep Linear Recurrent Units (LRUs), have yielded strong outcomes in language modeling tasks~\citep{Orvieto:23, De:24arxiv}. Hierarchically Gated Linear RNN (HGRN) models~\citep{Qin:23}, along with the subsequent HGRN2 variant~\citep{Qin:24arxiv}, have demonstrated similar promise. Notably, the RWKV architecture~\citep{Peng:23arxivshort,Peng:24arxivshort} stands out as an established RNN-based method in large language modeling, delivering results on par with those achieved by Transformers.

\paragraph{Gating.}
Gating, a fundamental concept introduced in LSTM, has been revisited and adapted in a variety of contemporary models. Numerous recent works, including HGRN~\citep{Qin:23}, HGRN2~\citep{Qin:24arxiv}, Gated Linear Attention (GLA)~\citep{Yang:23arxiv}, Gated State Space (GSS) models~\citep{Mehta:22}, Bidirectional Gated SSM (BiGS)~\citep{Wang:22}, Moving Average Equipped Gated Attention (MEGA)~\citep{Ma:22}, RWKV~\citep{Peng:23arxivshort}, and Mamba~\citep{Gu:24arxiv}, have incorporated gating mechanisms in novel ways.

\paragraph{Covariance Update Rule.} To improve storage capacity, we integrated a matrix-based memory into the mLSTM cell, governed by a covariance update rule. Several existing approaches also utilize similar update schemes, including Fast Weight Programmers \citep{Schmidhuber:92ncfastweights,Schlag:21}, RWKV-5 and RWKV-6~\citep{Peng:24arxivshort}, Retention~\citep{Sun:23arxiv}, Linear Transformer~\citep{Katharopoulos:20}, and HGRN2~\citep{Qin:24arxiv}.

\paragraph{Most Related.} The models most akin to xLSTM in terms of underlying principles are Retention~\citep{Sun:23arxiv}, RWKV~\citep{Peng:23arxivshort,Peng:24arxivshort}, and HGRN2~\citep{Qin:24arxiv}, all of which employ matrix-based memory and/or incorporate gating mechanisms. Nevertheless, these frameworks differ from the newly introduced sLSTM, as they lack support for memory mixing. The capability of memory mixing is crucial for addressing state tracking challenges, thereby granting LSTMs a greater degree of expressiveness compared to State Space Models (SSMs) and Transformers~\citep{Merrill:24,Deletang:23}. Effective state tracking is essential, for instance, when executing code or monitoring entities across extended narratives.

\paragraph{Residually Stacking Architectures.} Consistent with the majority of modern large-scale deep learning frameworks, xLSTM models are designed by arranging their fundamental components in a residual stacking fashion~\citep{Srivastava:15,He:16}.

This approach was pivotal in facilitating the success of deep convolutional neural networks~\citep{He:16} as well as the development of Transformer models~\citep{Vaswani:17}. Transformers, in particular, serve as the backbone for the current generation of Large Language Models (LLMs), such as GPT-3~\citep{Brown:20short}, ChatGPT~\citep{Schulman:22short}, GPT-4~\citep{Achiam:23gpt4short}, Megatron-LM~\citep{Shoeybi:19}, Gopher~\citep{Rae:21short}, ERNIE 3.0 Titan~\citep{Wang:21arxivshort}, GLaM~\citep{Du:21glamshort}, the Chinese M6~\citep{Lin:21}, the multilingual AlexaTM 20B~\citep{Soltan:22}, OPT~\citep{Zhang:22opt}, Chinchilla~\citep{Hoffmann:22short}, BLOOM~\citep{Scao:22arxivshort}, GLM-130B~\citep{Zeng:22arxivshort}, LaMDA~\citep{Thoppilan:22short}, PaLM~\citep{Chowdhery:22short}, Llama~\citep{Touvron:23arxiv}, and Gemini~\citep{GeminiTeam:23arxiv,Reid:24arxiv}.

\section{Experiments}
\label{sec:Exp}

This section presents an empirical analysis of xLSTM, positioning it relative to prevailing approaches in language modeling. Section~\ref{sec:ExpSyntethic} explores the distinct properties of xLSTM through a suite of synthetic benchmarks. Section~\ref{sec:ExpComparison} details a comparison of validation set perplexity across state-of-the-art language modeling techniques, each trained with 15B tokens from the SlimPajama dataset~\citep{Soboleva:23}; within this context, we also carry out ablation experiments targeting xLSTM. Additionally, we analyze how the methods scale, drawing on methodologies described by \citet{Kaplan:20} and \citet{Brown:20short}.

Subsequently, Section~\ref{sec:ExpLanguage} extends this investigation through comprehensive language modeling experiments. Here, we contrast xLSTM with the highest-performing models identified in Section~\ref{sec:ExpComparison}, using 300B SlimPajama~\citep{Soboleva:23} training tokens. Our analysis comprises four facets: (1) examining the ability of these methods to generalize to substantially longer input contexts; (2) measuring validation perplexity and evaluating outcomes on a suite of downstream tasks~\citep{Sutawika:24}; (3) assessing performance across 571 domains from the PALOMA language benchmark dataset~\citep{Magnusson:23arxivshort}; and (4) revisiting scaling dynamics in the context of a twenty-fold increase in training data.

\label{sec:xlstm_blocks_notation}
In our experiments, we denote xLSTM[$a$:$b$] to represent a configuration in which the ratio of mLSTM-based to sLSTM-based xLSTM blocks is $a/b$. As an illustration, xLSTM[7:1] signifies that among eight total blocks, seven utilize mLSTM while one utilizes sLSTM. When the overall number of blocks is set to 48, this yields 42 mLSTM-based blocks alongside 6 sLSTM-based blocks. Additionally, all experimental setups employ pre up-projection blocks for mLSTM and post up-projection blocks for sLSTM.

\subsection{Synthetic Tasks and Long Range Arena}
\label{sec:ExpSyntethic}
To begin, we evaluate the capability of xLSTM’s novel exponential gating mechanism with memory mixing on formal language benchmarks~\citep{Deletang:23}. Next, we examine how well xLSTM's matrix-based memory functions on the Multi-Query Associative Recall task~\citep{Arora:23arxiv}. Lastly, we measure xLSTM’s proficiency in handling extended sequences using the Long Range Arena benchmark~\citep{Tay:21}.

\paragraph{Evaluation of xLSTM's Exponential Gating with Memory Mixing.} We evaluate the novel exponential gating mechanism with memory mixing introduced in xLSTM, which is designed to facilitate state tracking capabilities~\citep{Merrill:24, Merrill:23}. Building on the formal language benchmarks from \citet{Deletang:23}, we modify and broaden these tasks to support training across multiple sequence lengths, thus allowing assessment of extrapolation abilities. Comprehensive descriptions of these tasks along with further results are provided in Appendix~\ref{sec:appExpSynthetic-formalLang}. In our comparisons, we benchmark xLSTM against a variety of architectures including Transformers, State Space Models, and Recurrent Neural Networks. We evaluate task accuracy specifically on those tokens pertinent to each respective formal language problem, scaling scores from 0 (chance performance) to 1 (perfect accuracy). Across all tasks, we compare the following 2-block model variants: xLSTM[0:1] (sLSTM only), xLSTM[1:0] (mLSTM only), xLSTM[1:1], as well as Llama, Mamba, RWKV, Retention, Hyena, standard LSTM, and LSTM embedded within Transformer blocks (LSTM (Block)). The resulting performance is displayed in Figure~\ref{fig:formal-main}. Notably, models such as Transformers or State Space Models that do not employ memory mixing—and therefore lack effective state tracking—are unable to resolve regular grammar challenges, such as the parity task. These observations corroborate prior evidence that both Transformers and State Space models possess inherently weaker representational capacity for such problems compared to RNNs~\citep{Merrill:24, Merrill:23, Deletang:23}.

\paragraph{Test of xLSTM's Memory Capacities on Associative Recall Tasks.} This experiment evaluates the enhanced matrix-based memory of xLSTM through its performance on the Multi-Query Associative Recall task~\citep{Arora:23arxiv}. In this setup, each sequence comprises randomly sampled key-value pairs from an extensive vocabulary, requiring the model to store and subsequently retrieve these associations. To make the task more demanding than in its original form, the number of key-value pairs per sequence is increased up to 256, and the input context is extended to a length of 2048, providing a more comprehensive examination of the memory capabilities across different models. We benchmark two-block variants of Llama, Mamba, RWKV-5, RWKV-6, xLSTM[1:1], and xLSTM[1:0] by measuring their accuracy in recalling the correct key-value pairs. Because Transformer architectures such as Llama possess memory that scales exponentially with their representation dimension~\citep{Ramsauer:21}, they are treated as the upper baseline for this benchmark. The results, depicted in Figure~\ref{fig:mqar-main}, show that xLSTM[1:1] achieves the highest accuracy among non-Transformer approaches, including configurations with fewer parameters. Notably, incorporating the sLSTM block does not impede memory performance; rather, it appears to amplify it, especially apparent in the scenario involving 256 key-value pairs, which represents the highest difficulty level assessed. Further detailed outcomes are provided in Appendix~\ref{sec:appExpSynthetic-mqar}, including extrapolation studies suggesting that the increased memory capacity of xLSTM supports generalization to sequence lengths that exceed those encountered during training.

\paragraph{Test of xLSTM's Long Context Capabilities on Long Range Arena.} To evaluate the ability of xLSTM to manage extended sequences and substantial contexts, we conduct comparisons with various approaches on the Long Range Arena~\citep{Tay:21}. Across all tasks, xLSTM achieves robust and reliable results, indicating that its architecture is highly effective in addressing the diverse challenges associated with long-context modeling.

Further information is available in Appendix~\ref{sec:appExpSynthetic-lra}.

\subsection{Method Comparison and Ablation Study}
\label{sec:ExpComparison}

In order to investigate the central question posed by our study—namely, the capabilities of our novel LSTM variants when scaled for language modelling—we conduct experiments training xLSTMs, Transformers, State Space Models, and additional approaches using 15B tokens sourced from SlimPajama under a uniform auto-regressive framework. These models are then evaluated against a validation dataset, and we carry out ablation studies specifically focused on the xLSTMs.

\paragraph{Comparing xLSTM to Other Methods.} To benchmark performance, we trained each model on a dataset of 15B tokens sampled from SlimPajama~\citep{Soboleva:23}, and assessed model quality using perplexity on a held-out validation set. The compared architectures include xLSTM (our proposed approach), GPT-3 (Transformer)~\citep{Brown:20short}, Llama (Transformer)~\citep{Touvron:23arxiv}, H3 (SSM)~\citep{Fu:23}, Mamba (SSM)~\citep{Gu:24arxiv}, RWKV-4, RWKV-5, and RWKV-6 (all RNNs)~\citep{Peng:23arxivshort, Peng:24arxivshort}, GLA (a linear Transformer variant)~\citep{Yang:23arxiv}, HGRN and HGRN2 (RNNs)~\citep{Qin:23, Qin:24arxiv}, as well as RetNet and Hyena (both linear Transformers)~\citep{Sun:23arxiv, Poli:23}, alongside xLSTM[1:0] and xLSTM[7:1] variants (refer to Section~\ref{sec:xlstm_blocks_notation}). All training runs used mixed precision, with the exception of RWKV-5, RWKV-6, GLA, and HGRN2, where mixed-precision training was manually managed instead of utilizing PyTorch's automated routines (see Appendix Section~\ref{sec:appGenTrainProc} for more information).

The experimental approaches are grouped as follows: (a) Transformers, (b) State Space Models (SSMs), and (c) Recurrent Neural Networks (RNNs), along with the class of linear Transformers. Notably, linear Transformers provide a linear alternative to the conventional attention mechanism found in standard Transformers. To ensure a fair comparison, all models are architecturally matched to a GPT-3 configuration containing approximately 350M parameters—this entails an embedding dimension of 1024 and 24 residual blocks. GPT-3 is the only model employing shared parameters between token and output embeddings, resulting in a reduced parameter count.

The results, presented in Table~\ref{tab:spaj15b_model_results}, demonstrate that xLSTM achieves superior perplexity compared to all existing methods tested. Further experimental details can be found in Appendix~\ref{sec:appExpComparison}. Additionally, Figure~\ref{fig:temp_sclaw15B} depicts model scaling trends in this setting, suggesting that the advantage of xLSTM persists, or potentially increases, as model size grows.

\paragraph{Ablation Studies.}
As shown in Table~\ref{tab:spaj15b_model_results} and Figure~\ref{fig:temp_sclaw15B}, xLSTM demonstrates strong performance on language modeling tasks after being trained on 15B tokens from SlimPajama. To systematically examine the individual contributions that differentiate xLSTM from a conventional LSTM, we incrementally transform a standard LSTM architecture toward an xLSTM. The process begins by embedding LSTM layers within a pre-LayerNorm residual structure. This configuration is then expanded to include a post up-projection block. In the final step, we incorporate exponential gating alongside matrix memory mechanisms.

Table~\ref{tab:ablstudies} (top) presents the findings.

The ablation experiments identify both the exponential gating and the matrix memory as key contributors to the observed performance gains. Given the centrality of gating in both RNNs and State Space Models, we further investigate alternative gating strategies. Results in Table~\ref{tab:ablstudies} (bottom) indicate that allowing each gate to be independently trainable and responsive to the input yields progressively better outcomes. Further analysis of specific backbone components can be found in Appendix~\ref{sec:appAblDetails}.

\subsection{xLSTM as Large Language Model}
\label{sec:ExpLanguage}

To conclude our investigation, we conduct extensive language modeling experiments to evaluate the efficacy of xLSTM in the context of large language models (LLMs). For this purpose, we scale up the training dataset, utilizing 300B tokens sourced from SlimPajama. Notably, this token count is consistent with that employed in prior works such as Mamba~\citep{Gu:24arxiv} and Griffin~\citep{De:24arxiv}.

We conduct a comparative analysis of xLSTM against RWKV-4, Llama, and Mamba, with each model representing a distinct methodological category outlined in Section~\ref{sec:ExpComparison}. RWKV-4 is chosen as the RNN baseline because satisfactory training precision configurations for RWKV-5, RWKV-6, and HGRN2 were only established subsequent to the commencement of the 300B token training runs (refer to Appendix~\ref{sec:appGenTrainProc}). Our evaluation encompasses multiple model scales (125M, 350M, 760M, 1.3B), and for each model, we analyze their ability to extrapolate across sequence lengths and examine their performance on validation datasets. Further, we measure efficacy on various downstream benchmarks, test language modeling accuracy over 571 text domains from the PALOMA benchmark, and explore how each model adheres to scaling laws.

\paragraph{Sequence Length Extrapolation.}
\label{sec:spaj300B_ctx_extrapolation}
We begin by evaluating the ability of large-scale models, specifically those with 1.3B parameters—xLSTM, RWKV-4, Llama, and Mamba—to generalize to longer sequence lengths. Each model is initially trained with a context length of 2048, after which performance is assessed for increasingly extended contexts, up to 16384 tokens. Figure~\ref{fig:spaj300B_seq_extrapolation} presents these findings. Notably, the xLSTM model stands out for preserving low perplexity even as the context length increases, distinguishing itself from the other architectures.

\paragraph{Validation Perplexity and Downstream Tasks.}
\label{sec:spaj_downstream_eval}
Next, across each model size, we assess the xLSTM, RWKV-4, Llama, and Mamba architectures using the SlimPajama validation set for next token prediction accuracy as well as on various downstream benchmarks that test common sense reasoning capabilities.

The third column in Table~\ref{tab:lmeval} displays the validation set perplexity values obtained by the different approaches. Across every model size, xLSTM[1:0] and xLSTM[7:1] yield the lowest validation set perplexities, indicating superior performance in this metric. The subsequent columns in Table~\ref{tab:lmeval} summarize results for various downstream tasks. In nearly all evaluated tasks and for all model scales, xLSTM consistently outperforms competing methods, with the exception of the ARC task where Mamba occasionally achieves top performance. Additional information can be found in Appendix~\ref{sec:appExpLanguage}.

\paragraph{Performance on PALOMA Language Tasks.} As a third evaluation, we assess the next token prediction capabilities of xLSTM, RWKV-4, Llama, and Mamba across all model sizes using the PALOMA language benchmarks~\citep{Magnusson:23arxivshort}. Performance is evaluated using perplexity for next token prediction over a set of 571 text domains, encompassing sources such as nytimes.com and the Reddit forum r/depression. Table~\ref{tab:ppleval} summarizes token prediction perplexity, categorizing results into language modeling tasks (the first seven columns) and more specific domain evaluations (the last five columns). On these benchmarks, xLSTM[1:0] outperforms xLSTM[7:1] with respect to language task performance.

In comparison to Mamba, xLSTM[1:0] achieves lower perplexity in 568 out of 571 text domains (99.5\%). When evaluated against Llama, it demonstrates lower perplexity in 486 of 571 domains (85.1\%), and versus RWKV-4, xLSTM[1:0] outperforms in 570 out of 571 cases (99.8\%). Further information can be found in Appendix~\ref{sec:appExpLanguage}.

\paragraph{Scaling Laws.} Next, we examine the power-law scaling properties, which enable predictions about model performance at increased parameter counts~\citep{Kaplan:20,Brown:20short}. The scaling trends are illustrated in Figure~\ref{fig:temp_sclaw300B}. Although all the models exhibit analogous scaling patterns, they differ in their baseline performance levels. RWKV-4 shows the least favorable results, succeeded by Llama and Mamba. Notably, xLSTM outperforms Mamba, maintaining approximately the same advantage over Mamba as Mamba has over Llama. These observed scaling dynamics suggest that as model sizes expand, xLSTM is likely to maintain its performance edge over both Transformer-based and State-Space models.

\textbf{Generation Times and Maximal Throughput.} We report the results for text generation latency in Figure~\ref{fig:inference_time_generation} and assess maximal decoding throughput in Figure~\ref{fig:inference_time_decoding_throughput} (left) across our xLSTM variants at the 1.3B parameter scale. Our evaluation includes implementations of Mamba, Llama, and RWKV with comparable model sizes sourced from HuggingFace, where the Llama model incorporates a static key-value cache. It is notable that, at the time of these evaluations, neither complete cache compilation for the Transformer architecture nor successful compilation of Mamba models using \texttt{torch.compile} were feasible. All models are evaluated using a batch size of 1 with a pre-fill value of 16 during text generation, which is optimal for Transformer-based models.

The data in Figure~\ref{fig:inference_time_generation} demonstrate that xLSTM, along with other recurrent frameworks such as Mamba and RWKV-4, scales linearly with sequence length, whereas Llama exhibits quadratic time complexity. For maximal decoding throughput, we consider variations in batch size and prefill parameters for the Llama baseline. As illustrated in Figure~\ref{fig:inference_time_decoding_throughput} (right), the constant memory requirement of xLSTM enables support for substantially larger batch sizes relative to Llama, resulting in superior throughput performance.

\section{Limitations}

(i) Unlike mLSTM, the sLSTM's approach to memory mixing does not support operations that can be easily parallelized, restricting the possibility for rapid parallel computation. Despite this limitation, we engineered an efficient CUDA kernel for sLSTM; this implementation currently performs at less than twice the runtime of our parallel mLSTM variant. (ii) The CUDA kernels developed for mLSTM have not undergone optimization, resulting in execution that is approximately four times slower compared to FlashAttention or the scan mechanism implemented in Mamba. Optimized CUDA kernels, modeled after FlashAttention, could considerably improve the processing speed. (iii) The use of matrix memory in mLSTM introduces significant computational complexity because it requires manipulation of $d \times d$ matrices. Nonetheless, memory updates and access are non-parametric and can be performed in parallel using common matrix operations, which means the increased wall clock time associated with this complex memory structure is relatively insignificant. (iv) Proper selection of initialization values for the forget gates is essential. (v) As the matrix memory in mLSTM does not depend on the input sequence length, extending the sequence may saturate the memory when dealing with longer contexts. However, for context lengths up to 16k, this has not posed a practical limitation, as discussed in Section~\ref{sec:spaj300B_ctx_extrapolation}. (vi) Given the high computational demands inherent in large language model training, we did not conduct exhaustive architectural or hyperparameter tuning, particularly for the larger xLSTM models. Achieving peak performance from xLSTM is likely contingent on a more comprehensive optimization effort.

\section{Conclusion}
We have made progress toward answering our central question: To what extent can LSTM be scaled, in terms of parameter count, for language modeling? Our results demonstrate that scaling LSTM to billions of parameters yields performance that matches, at a minimum, the capabilities of contemporary architectures such as Transformers and State Space Models. Through introducing the xLSTM variant—incorporating exponential gating with memory mixing alongside a revised memory structure—we observe that xLSTM achieves strong results in language modeling, standing on par with, or surpassing, leading approaches like Transformers and State Space Models. Observations of scaling behaviors suggest that, as xLSTM models grow larger, they may emerge as formidable alternatives to the predominant Transformer-based Large Language Models. Additionally, xLSTM holds promise for making significant advances in domains beyond language, including Reinforcement Learning, Time Series Forecasting, and physical systems modeling.

\end{document}